\chapter{Methodology}

\section{User Study}

\subsection*{Study Design}
\begin{itemize}
    \item \textbf{Type:} Within-subjects experiment with counterbalanced order using a Balanced Latin Square design
    \item \textbf{Sample:} N = 73 university students (novice users, no prior GDPR experience)
    \item \textbf{Duration:} Approximately 90 minutes per participant
    \item \textbf{Counterbalancing:}
    \begin{itemize}
        \item Balanced Latin Square with 6 groups (reflecting 3 AI modes tested in randomized sequences)
        \item Each participant completes all three modes in a predetermined counterbalanced order
        \item \textit{Rationale:} Controls for order effects, learning effects, and fatigue
    \end{itemize}
\end{itemize}

\subsection*{Task Scenario}
\begin{itemize}
    \item \textbf{Consistency:} Single standardized scenario used across all participants and modes
    \item \textbf{Content:} A relatable university-based data processing activity covering all GDPR Article 30 fields
    \item \textit{Rationale:} Ensures comparability across conditions while remaining accessible to participants without legal expertise
    \item \textbf{LLM Model:} Mistral (consistent across all modes)
\end{itemize}

\subsection*{Participants}
\begin{itemize}
    \item \textbf{Recruitment:} N = 73 university students
    \item \textbf{Screening Criteria:}
    \begin{itemize}
        \item Minimal to no GDPR knowledge (expected and required)
        \item Prior AI tool usage experience (documented)
        \item General technology comfort level (assessed)
        \item No professional legal or compliance background
    \end{itemize}
    \item \textbf{Onboarding:} Participants receive a brief standardized introduction to GDPR and ROPA concepts (e.g., short instructional video or text) before beginning the tasks
    \item \textbf{Reference Material:} Participants are shown an example of a high-quality ROPA document prior to starting
\end{itemize}

\subsection*{Data Collection}
\begin{itemize}
    \item \textbf{Quantitative Measures (collected after each mode):}
    \begin{itemize}
        \item System Usability Scale (SUS): 10 standardized questions
        \item Raw Task Load Index (R-TLX): 6 cognitive load dimensions
        \item LLM-Specific Confidence Questions (7-point Likert scale):
        \begin{enumerate}
            \item \textit{Perceived Usefulness:} ``The LLM was helpful in creating this document.''
            \item \textit{Efficiency:} ``This mode let me complete the ROPA faster than I could on my own.''
            \item \textit{Transparency:} ``I understood why the AI suggested certain GDPR fields.''
            \item \textit{Reliance Comfort:} ``I felt comfortable relying on the AI's suggested categories.''
            \item \textit{Task-Specific Confidence:} ``I am confident the AI identified the appropriate legal basis.''
            \item \textit{Intention to Use:} ``I would use this mode again for similar work.''
            \item \textit{Professional Applicability:} ``If I were responsible for GDPR in a real company, I would use this specific mode to generate my documentation.''
        \end{enumerate}
        \item Free-text qualitative feedback (collected after each mode)
        \item Automatic completion time tracking (per mode)
    \end{itemize}
    \item \textbf{Post-Experiment Ranking:}
    \begin{itemize}
        \item Participants rank all three modes by overall preference after completing all conditions
    \end{itemize}
\end{itemize}

\subsection*{Quality Assessment}
\begin{itemize}
    \item \textbf{BERTScore Evaluation:}
    \begin{itemize}
        \item Compare all documents (73 participants × 3 modes) against an expert-created reference (ECR)
        \item Aggregate scores by mode to establish baseline quality metrics
        \item Perform pairwise comparisons between modes and against the ECR for comprehensive evaluation
    \end{itemize}
    \item \textbf{LLM-as-a-Judge Validation:}
    \begin{itemize}
        \item Evaluate a representative subset of mode-generated documents using Claude
        \item Cross-validate BERTScore findings with LLM-based quality assessment
        \item \textit{Rationale:} Ensures BERTScore captures meaningful quality differences beyond lexical similarity
    \end{itemize}
    \item \textbf{Confidence-Quality Gap Analysis:}
    \begin{itemize}
        \item Correlate self-reported confidence scores with actual document quality (BERTScore)
        \item Identify potential over-confidence or under-confidence patterns across modes
    \end{itemize}
\end{itemize}

\subsection*{Statistical Analysis}
\begin{itemize}
    \item \textbf{RQ1 (User Experience \& Preference):}
    \begin{itemize}
        \item Repeated measures ANOVA on SUS scores, R-TLX dimensions, and LLM confidence questions
        \item Post-hoc pairwise comparisons (with Bonferroni correction) to identify significant differences between modes
        \item Thematic analysis of free-text qualitative feedback and post-experiment rankings
    \end{itemize}
    \item \textbf{RQ2 (Document Quality \& Confidence Calibration):}
    \begin{itemize}
        \item Repeated measures ANOVA on BERTScore results across modes
        \item Correlation analysis between self-reported confidence and actual quality (BERTScore)
    \end{itemize}
\end{itemize}